\chapter[PID Modules and Canonical Forms]{Modules over PIDs and Canonical Forms}\label{chap:pid-modules-canonical-forms}

Chapter~5 showed that diagonalization is possible exactly when an operator has
enough eigenvectors. This chapter asks what remains when it does not. The
answer comes from viewing a linear operator as multiplication by \(x\), and
from the remarkable fact that finitely generated modules over a principal
ideal domain have a complete decomposition theory. The same theorem therefore
classifies finitely generated abelian groups and produces rational and Jordan
canonical forms.

The first half of the chapter deliberately uses three different viewpoints.
Compatible ambient coordinates construct the decomposition, primary layers
recognize its pieces intrinsically, and Smith reduction expresses the result in computational
matrix language.  Keeping those tasks separate is the key to the proof
architecture.

Through the Smith normal form section, \(R\) is a principal ideal domain (PID).
A PID need not be a Euclidean domain. We use only principal ideals, gcds and
B\'ezout identities; no division-with-remainder argument is used over an
arbitrary PID. The polynomial ring \(F[x]\), used later, is Euclidean, so its
ordinary polynomial division algorithm is available there.

\section{Submodules of Free Modules over a PID}\label{sec:pid-submodules-free}
Over a field, every subspace has a basis. The analogous statement fails over
general rings, but a PID retains just enough divisibility control for finite
free modules to behave similarly.

\begin{theorem}[Submodules of finite free modules over a PID]
\label{thm:pid-submodules-compatible-bases}
Let \(R\) be a PID, let \(F\) be a free \(R\)-module of finite rank \(n\),
and let \(N\subseteq F\) be a submodule.
\begin{enumerate}
\item The module \(N\) is free of rank at most \(n\).
\item There exist a basis \(y_1,\ldots,y_n\) of \(F\), an integer
      \(m\leq n\), and nonzero elements \(a_1,\ldots,a_m\in R\) such that
\[
a_1\mid a_2\mid\cdots\mid a_m
\qquad\text{and}\qquad
N=Ra_1y_1\oplus\cdots\oplus Ra_my_m.
\]
\end{enumerate}
Part~(1) is the freeness conclusion. Part~(2) strengthens it by choosing
compatible coordinates in the ambient module; it makes no claim about an
arbitrarily prescribed basis of \(F\).
\end{theorem}
\begin{proof}
\emph{Proof of (1).}
Choose a basis of \(F\) and identify it with \(R^n\). We argue by induction
on \(n\). The assertion is clear for \(n=0\). Let
\[
\pi:R^n\longrightarrow R
\]
be projection onto the first coordinate, and set
\[
K=\ker(\pi|_N)=N\cap(0\oplus R^{n-1}).
\]
The image \(\pi(N)\) is an ideal, say \((d)\).

If \(\pi(N)=0\), then \(N\subseteq0\oplus R^{n-1}\). Induction makes \(N\)
free of rank at most \(n-1\), and hence at most \(n\).

Suppose instead that \(\pi(N)=(d)\ne0\). Choose \(x\in N\) with
\(\pi(x)=d\). For every \(y\in N\), there is \(r\in R\) such that
\(\pi(y)=rd\). Therefore
\[
\pi(y-rx)=rd-rd=0,
\]
so \(y-rx\in K\), and
\[
y=rx+(y-rx)\in Rx+K.
\]
This proves \(N=Rx+K\). The sum is direct: if \(ax+k=0\) with \(k\in K\),
then applying \(\pi\) gives
\[
ad=\pi(ax+k)=0.
\]
Because \(R\) is a domain and \(d\ne0\), we have \(a=0\), and then \(k=0\).

By induction, choose a basis \(k_1,\ldots,k_s\) of \(K\), with
\(s\leq n-1\). The decomposition above shows that
\(x,k_1,\ldots,k_s\) spans \(N\), and the directness calculation together
with independence of the \(k_i\) shows that this list is independent. Hence
it is a basis of \(N\), and
\[
\operatorname{rank}N=1+s\leq n.
\]

\medskip
\noindent\emph{Proof idea for (2).}
We find one invariant-factor direction at a time.  Among all linear
functionals \(\phi:F\to R\), choose one whose image \(\phi(N)\) is maximal
as an ideal.  Its generator becomes the first invariant factor.  We show that
a corresponding element of \(N\) is a scalar multiple of a primitive basis
vector, split this direction from both \(F\) and \(N\), and repeat inside a
free module of rank one smaller:
\[
\begin{gathered}
\text{choose the first coefficient}\longrightarrow
\text{find a primitive direction},\\
\text{split it off}\longrightarrow\text{induct}.
\end{gathered}
\]

\medskip
\noindent\emph{Proof of (2).}
We induct on \(n\).  The result is clear if \(n=0\), and if \(N=0\) we may
take \(m=0\) and any basis of \(F\).  Suppose therefore that \(n>0\) and
\(N\ne0\).

\emph{Stage 1: choose a maximal image ideal.}
For every \(\phi\in\operatorname{Hom}_R(F,R)\), the set \(\phi(N)\) is an
ideal.  A PID is Noetherian because all its ideals are principal, so the
maximal-member characterization of Noetherian rings proved in Chapter~4
gives a maximal member of the partially ordered set
\[
\Sigma=\{\phi(N):\phi\in\operatorname{Hom}_R(F,R)\},
\qquad\text{ordered by set inclusion}.
\]
Choose \(\psi:F\to R\) with \(\psi(N)=(a_1)\) maximal.  This ideal is
nonzero. Indeed, choose \(0\ne z\in N\) and a basis
\(x_1,\ldots,x_n\) of \(F\). Write
\[
z=r_1x_1+\cdots+r_nx_n.
\]
Some \(r_i\ne0\). Define the projection onto the \(i\)-th coordinate
relative to this basis by
\[
\pi_i:F\longrightarrow R,
\qquad
\pi_i\left(\sum_{j=1}^n r_jx_j\right)=r_i.
\]
Then \(\pi_i(z)=r_i\ne0\), so \(\pi_i(N)\ne0\) is a nonzero member of
\(\Sigma\). Recall the divisibility reversal
\[
(a)\subseteq(b)\quad\Longleftrightarrow\quad b\mid a.
\]
Thus enlarging an image ideal means finding a more fundamental divisor.

\emph{Stage 2: force divisibility of every scalar measurement.}
Choose \(z_1\in N\) with \(\psi(z_1)=a_1\).  We claim that
\[
a_1\mid\phi(z_1)
\qquad\text{for every }\phi:F\to R.
\]
If one scalar measurement of \(z_1\) were not divisible by \(a_1\), the
B\'ezout calculation below would produce a strictly larger image ideal.

Fix an arbitrary \(\phi:F\to R\). Let \((a_1,\phi(z_1))=(g)\), and choose
\(r,s\in R\) such that
\(g=ra_1+s\phi(z_1)\). Define \(\nu=r\psi+s\phi\). Then
\[
\nu(z_1)=r\psi(z_1)+s\phi(z_1)
=ra_1+s\phi(z_1)=g.
\]
Since \(z_1\in N\), this gives \(g\in\nu(N)\), and hence
\((g)\subseteq\nu(N)\). Also \(g\mid a_1\), so principal-ideal inclusion
reverses divisibility and gives \((a_1)\subseteq(g)\). Therefore
\[
(a_1)\subseteq(g)\subseteq\nu(N).
\]
The map \(\nu\) belongs to \(\operatorname{Hom}_R(F,R)\), so \(\nu(N)\)
is a member of \(\Sigma\). Since \((a_1)=\psi(N)\) was chosen maximal in
\(\Sigma\), the inclusion \((a_1)\subseteq\nu(N)\) forces
\(\nu(N)=(a_1)\). Consequently \((a_1)=(g)\), so \(g\) is associate to
\(a_1\). Since \(g\mid\phi(z_1)\), it follows that
\(a_1\mid\phi(z_1)\).

\emph{Stage 3: identify a primitive direction.}
Choose any basis \(x_1,\ldots,x_n\) of \(F\), and define the projection onto
the \(i\)-th coordinate relative to this basis by
\[
\pi_i:F\longrightarrow R,
\qquad
\pi_i\left(\sum_{j=1}^n r_jx_j\right)=r_i.
\]
Since the preceding step applies to every \(\pi_i\), write
\(\pi_i(z_1)=a_1b_i\) and set
\[
y_1=\sum_{i=1}^n b_ix_i.
\]
The coordinate expansion of \(z_1\) now gives
\[
z_1=\sum_i\pi_i(z_1)x_i
=\sum_i a_1b_ix_i=a_1y_1.
\]
Moreover,
\[
a_1=\psi(z_1)=\psi(a_1y_1)=a_1\psi(y_1).
\]
Because \(R\) is a domain and \(a_1\ne0\), we have \(\psi(y_1)=1\).
Thus \(y_1\) has an \(R\)-linear functional taking it to \(1\). This is the
key primitive property: it makes \(Ry_1\) split off from \(F\).

\emph{Stage 4: split \(F\) and \(N\) simultaneously.}
For every \(x\in F\),
\[
x=\psi(x)y_1+\bigl(x-\psi(x)y_1\bigr).
\]
Indeed,
\[
\psi\bigl(x-\psi(x)y_1\bigr)
=\psi(x)-\psi(x)\psi(y_1)=0,
\]
so the second term lies in \(\ker\psi\), and hence
\(F=Ry_1+\ker\psi\). If \(ry_1\in\ker\psi\), then
\[
0=\psi(ry_1)=r\psi(y_1)=r.
\]
Thus \(Ry_1\cap\ker\psi=0\), and therefore
\[
F=Ry_1\oplus\ker\psi.
\]
If \(z\in N\), then \(\psi(z)=ra_1\) for some \(r\in R\), so
\(z-ra_1y_1\in N\), and
\[
\psi(z-ra_1y_1)=ra_1-ra_1\psi(y_1)=0.
\]
Therefore
\[
z=ra_1y_1+(z-ra_1y_1),
\]
with the second summand in \(N\cap\ker\psi\). This proves
\(N=Ra_1y_1+(N\cap\ker\psi)\). If
\(ra_1y_1\in N\cap\ker\psi\), then
\[
0=\psi(ra_1y_1)=ra_1.
\]
Since \(R\) is a domain and \(a_1\ne0\), we have \(r=0\). Hence the sum is
direct:
\[
N=Ra_1y_1\oplus(N\cap\ker\psi).
\]
One relation direction has now been completely separated from all the others.

\emph{Stage 5: induct and recover the divisibility chain.}
The module \(\ker\psi\) is a submodule of \(F\), hence free by part~(1).
Since \(\operatorname{rank}F=n\), \(\operatorname{rank}(Ry_1)=1\), and
\(F=Ry_1\oplus\ker\psi\),
Proposition~\ref{prop:finite-direct-sum-rank} gives
\[
n=1+\operatorname{rank}(\ker\psi).
\]
Thus \(\operatorname{rank}(\ker\psi)=n-1\), which justifies the induction
parameter. Apply induction to
\(N\cap\ker\psi\subseteq\ker\psi\).  We obtain basis vectors
\(y_2,\ldots,y_n\) and nonzero elements \(a_2,\ldots,a_m\) with
\[
a_2\mid\cdots\mid a_m,
\qquad
N\cap\ker\psi=Ra_2y_2\oplus\cdots\oplus Ra_my_m.
\]
If \(m\geq2\), define \(\phi:F\to R\) on the new basis by
\[
\phi(y_1)=1,\qquad \phi(y_2)=1,\qquad
\phi(y_i)=0\quad(i\geq3).
\]
Then
\[
\phi(a_1y_1)=a_1,
\qquad
\phi(a_2y_2)=a_2.
\]
Hence \(a_1,a_2\in\phi(N)\), and in particular
\((a_1)\subseteq\phi(N)\). Because
\(\phi\in\operatorname{Hom}_R(F,R)\), the ideal \(\phi(N)\) is a member of
\(\Sigma\). Maximality of \((a_1)\) therefore gives
\(\phi(N)=(a_1)\).
Therefore \(a_2\in(a_1)\), or \(a_1\mid a_2\), as required.
\end{proof}

\begin{corollary}[Rank--nullity for finite free modules over a PID]
\label{cor:pid-rank-nullity}
Let \(R\) be a PID and let \(f:F\to G\) be an \(R\)-linear map between
finite free \(R\)-modules. Then \(\ker f\) and \(\operatorname{im}f\) are
finite free, the short exact sequence
\[
0\longrightarrow\ker f\longrightarrow F
\xrightarrow{f}\operatorname{im}f\longrightarrow0
\]
splits, and
\[
F\cong\ker f\oplus\operatorname{im}f,
\qquad
\operatorname{rank}F
=\operatorname{rank}(\ker f)+\operatorname{rank}(\operatorname{im}f).
\]
More precisely, after choosing a section
\(s:\operatorname{im}f\to F\), the literal internal decomposition is
\(F=\ker f\oplus s(\operatorname{im}f)\), with
\(s(\operatorname{im}f)\cong\operatorname{im}f\).
\end{corollary}
\begin{proof}
The inclusions \(\ker f\leq F\) and \(\operatorname{im}f\leq G\), together
with Theorem~\ref{thm:pid-submodules-compatible-bases}(1), show that both
modules are finite free. In particular \(\operatorname{im}f\) is projective,
so Theorem~\ref{thm:projective-characterizations} splits the displayed short
exact sequence. The Splitting Lemma gives the stated internal decomposition,
and Proposition~\ref{prop:finite-direct-sum-rank} gives the rank formula.
\end{proof}

One should not conclude similarly that
\(G\cong\operatorname{im}f\oplus\operatorname{coker}f\): although
\(2\mathbb Z\leq\mathbb Z\) is free, it is not a direct summand.

\section{Structure Theorem: Existence in Invariant-Factor Form}
\label{sec:pid-structure-existence}
A finitely generated module is built from generators subject to relations.
The compatible-basis conclusion of Theorem~\ref{thm:pid-submodules-compatible-bases}
separates those relations inside a finite free module; taking the quotient
then reads off the answer coordinate by coordinate.

\begin{theorem}[Structure theorem: existence]\label{thm:pid-structure-existence}
Let \(M\) be a finitely generated module over a PID \(R\).
\begin{enumerate}
\item There exist an integer \(r\geq0\) and nonzero nonunits
\(a_1,\ldots,a_m\) such that
\[
a_1\mid a_2\mid\cdots\mid a_m
\quad\text{and}\quad
M\cong R/(a_1)\oplus\cdots\oplus R/(a_m)\oplus R^r.
\]
\item The module \(M\) is free if and only if it is torsion-free.
\item In the decomposition above,
\[
T(M)\cong R/(a_1)\oplus\cdots\oplus R/(a_m),
\qquad M/T(M)\cong R^r.
\]
\end{enumerate}
\end{theorem}
\begin{proofidea}
Choose generators of \(M\) and realize it as a quotient \(F/N\) of a finite
free module. Apply Theorem~\ref{thm:pid-submodules-compatible-bases}(2) to
choose compatible coordinates for \(N\subseteq F\). The quotient becomes
coordinatewise, producing the invariant factors. Once the free and torsion
summands are visible, parts (2) and (3) follow immediately.
\end{proofidea}
\begin{proof}
Choose a surjection \(F=R^n\twoheadrightarrow M\) and let \(N\) be its
kernel. By Theorem~\ref{thm:pid-submodules-compatible-bases}(2), there are a basis
\(y_1,\ldots,y_n\) of \(F\) and nonzero \(a_1\mid\cdots\mid a_s\) such that
\[
N=Ra_1y_1\oplus\cdots\oplus Ra_sy_s.
\]
Define
\[
\Phi:F\longrightarrow
R/(a_1)\oplus\cdots\oplus R/(a_s)\oplus R^{n-s}
\]
by
\[
\Phi\left(\sum_{i=1}^n r_iy_i\right)
=\bigl(r_1+(a_1),\ldots,r_s+(a_s),r_{s+1},\ldots,r_n\bigr).
\]
Coordinatewise addition and scalar multiplication show directly that
\(\Phi\) is \(R\)-linear. It is surjective: choose representatives in \(R\)
for the first \(s\) coordinates of any target tuple and use its remaining
coordinates directly. Moreover,
\[
\Phi\left(\sum_i r_iy_i\right)=0
\]
if and only if \(r_i\in(a_i)\) for \(1\leq i\leq s\) and \(r_i=0\) for
\(i>s\). Thus \(r_i=a_ic_i\) for \(i\leq s\), and
\[
\ker\Phi=Ra_1y_1\oplus\cdots\oplus Ra_sy_s=N.
\]
The First Isomorphism Theorem therefore gives
\[
M\cong F/N
\cong R/(a_1)\oplus\cdots\oplus R/(a_s)\oplus R^{n-s}.
\]
If some \(a_i\) is a unit, then \(R/(a_i)=0\); omit every such zero summand.
Relabel the remaining divisibility chain as \(a_1,\ldots,a_m\).  Its cyclic
factors are nonzero and the free rank is \(r=n-s\), which proves part (1).

If \(M\) is free, it is torsion-free because \(R\) is a domain.  Conversely,
each nonzero summand \(R/(a_i)\) in part (1) is torsion, since
\[
a_i\bigl(1+(a_i)\bigr)=0.
\]
Thus a torsion-free \(M\) has no cyclic torsion summands and is free, proving
part (2).

Finally, set \(C=R/(a_1)\oplus\cdots\oplus R/(a_m)\). If \(m>0\), the
divisibility \(a_i\mid a_m\) makes \(a_m\) annihilate every cyclic summand,
hence every element of \(C\); if \(m=0\), then \(C=0\). Thus
\(C\subseteq T(M)\). Conversely, write an element of
\(M\cong C\oplus R^r\) as \(x=c+u\). If \(x\) is torsion, some
\(0\ne d\in R\) satisfies \(d(c+u)=0\). Its free coordinate gives
\(du=0\). Since \(R^r\) is torsion-free, \(u=0\), so every torsion element
lies in \(C\). Therefore \(T(M)=C\), and quotienting gives
\(M/T(M)\cong R^r\). This proves part (3).
\end{proof}

\section{Elementary Divisors and Primary Decomposition}
Invariant factors group the torsion into a divisibility chain.  Elementary
divisors reorganize the same module in the opposite direction: they separate
the independent prime-power phenomena.  This is still an existence argument;
intrinsic uniqueness comes afterward.

\begin{proposition}[Elementary-divisor existence]
The torsion part of a finitely generated PID module is a direct sum of modules
\(R/(p^e)\), where \(p\) is a prime element and \(e\geq1\).
\end{proposition}
\begin{proof}
Every PID is a UFD by Theorem~\ref{thm:pid-ufd}. Factor an invariant factor as
\[
d=u\prod_i p_i^{e_i},
\]
with distinct \(p_i\). The ideals
\((p_i^{e_i})\) are pairwise comaximal: a common divisor of two distinct
prime powers is a unit, so a B\'ezout identity gives their sum as \(R\).
The Chinese Remainder Theorem therefore gives
\[R/(d)\cong\bigoplus_iR/(p_i^{e_i}).\]
Apply this to every invariant factor.
\end{proof}

Fix one representative from each associate class of prime elements of \(R\).
When \(R=F[x]\), we take the monic irreducible polynomials as representatives.
For such a representative \(p\), define the \(p\)-primary component of \(M\) by
\[
M_p:=M[p^\infty]
=\{m\in M:p^em=0\text{ for some }e\geq1\}.
\]
Grouping the elementary-divisor summands by these representatives gives
submodules \(E_p\) and a decomposition
\[
T(M)=\bigoplus_p E_p.
\]
Each \(E_p\) is killed by a power of \(p\), so
\(E_p\subseteq M_p\). Conversely, suppose
\(m\in M_p\) and write \(m=\sum_qm_q\), with \(m_q\in E_q\).
For \(q\ne p\), some powers \(p^e\) and \(q^f\) kill \(m_q\). Since
distinct prime representatives are relatively prime, B\'ezout gives
\(up^e+vq^f=1\), and therefore
\(m_q=up^em_q+vq^fm_q=0\). Thus \(m=m_p\in E_p\), proving
\(E_p=M_p\). Hence the intrinsic primary decomposition is
\[
T(M)=\bigoplus_p M_p,
\]
where \(p\) runs over the fixed representatives and only finitely many
summands are nonzero. For example, the CRT
decomposition \(\mathbb Z/12\cong\mathbb Z/4\oplus\mathbb Z/3\) separates
the \(2\)-primary and \(3\)-primary behavior.

\section{Why Is the Decomposition Unique?}
Existence asked a constructive question: how can we choose coordinates that
split a module into standard pieces?  Uniqueness asks a different question:
how can those pieces be recognized from \(M\) without referring to chosen
coordinates?

The free part is immediate from Theorem~\ref{thm:pid-structure-existence}(3).
The torsion submodule \(T(M)\) is intrinsic, so
\[
M/T(M)\cong R^r.
\]
Invariant basis number from Chapter~5 therefore makes \(r\) intrinsic.  The
real problem is to recover the torsion blocks.

There is a useful first clue. Suppose a nonzero torsion module \(M\) has an
invariant-factor decomposition
\[
M\cong R/(a_1)\oplus\cdots\oplus R/(a_k),
\qquad a_1\mid\cdots\mid a_k.
\]
The annihilator of a direct sum is the intersection of the annihilators of
its summands: a scalar annihilates the direct sum exactly when it annihilates
every coordinate. The divisibility chain therefore gives
\[
\operatorname{Ann}(M)=(a_1)\cap\cdots\cap(a_k)=(a_k).
\]
Thus the largest invariant factor is visible intrinsically.  It is tempting
to remove a displayed copy of \(R/(a_k)\) and repeat.  But an abstract
isomorphism of direct sums need not identify the last displayed summand on
one side with the last displayed summand on the other.  Cancellation cannot
be assumed silently.  We instead seek stronger intrinsic invariants that
recover every block at once.

Different primes describe independent torsion phenomena, so the intrinsic
submodules \(M_p\) tell us to fix \(p\) and study \(M_p\), solving uniqueness
one prime at a time.

\section{Primary Layers and Uniqueness}\label{sec:primary-layers}
Let \(P\) be a finitely generated \(p\)-primary module.  Elementary-divisor
existence shows that one common power of \(p\) annihilates \(P\), so its powers
of \(p\) form a finite filtration
\[
P\supseteq pP\supseteq p^2P\supseteq\cdots\supseteq0.
\]
What disappears between \(p^{j-1}P\) and \(p^jP\)?  The quotient
\(p^{j-1}P/p^jP\) measures exactly that loss.  Multiplication by \(p\) kills
the quotient, so it is naturally a vector space over \(R/(p)\).  This
motivates the intrinsic layer dimensions
\[
a_j(P)=\dim_{R/(p)}\bigl(p^{j-1}P/p^jP\bigr)
\qquad(j\geq1).
\]
At the first layer, every cyclic \(p\)-power block is still present, so
\(a_1(P)\) counts all blocks. At the second layer only blocks of exponent at
least two survive, so \(a_2(P)\) counts those blocks; at the third layer only
blocks of exponent at least three survive, and so on. The successive layers
therefore record the block heights without choosing a decomposition.

First consider one block \(C=R/(p^4)\).  Its filtration is
\[
C\supseteq pC\supseteq p^2C\supseteq p^3C\supseteq p^4C=0,
\]
and every nonzero successive quotient is one-dimensional over \(R/(p)\).
Its layer sequence is therefore \(1,1,1,1,0,\ldots\).  Similarly,
\(R/(p^2)\) contributes \(1,1,0,\ldots\), while \(R/(p)\) contributes
\(1,0,\ldots\).

For the multi-block example
\[
P=R/(p^4)\oplus R/(p^2)\oplus R/(p^2),
\]
the Young-diagram bookkeeping appears numerically as
\[
\begin{array}{c|ccccc}
 &a_1&a_2&a_3&a_4&a_5\\ \hline
R/(p^4)&1&1&1&1&0\\
R/(p^2)&1&1&0&0&0\\
R/(p^2)&1&1&0&0&0\\ \hline
P&3&3&1&1&0.
\end{array}
\]
Representing each cyclic summand by a column whose height is its exponent,
\(a_j\) counts columns reaching height \(j\), while \(a_j-a_{j+1}\) counts
columns ending there.  The picture reveals the recovery mechanism before the
formal proof.

\begin{lemma}[Primary layers]\label{lem:primary-layers}
If
\[
P\cong\bigoplus_iR/(p^{e_i}),
\]
then
\[
a_j(P)=\#\{i:e_i\geq j\}.
\]
\end{lemma}
\begin{proof}
We separate the four mechanisms in the formula.

\emph{Vector-space structure.}
Because \(p\) is a nonzero prime element of the PID \(R\),
Proposition~\ref{prop:pid-nonzero-prime-maximal} shows that \(R/(p)\) is a
field.
Define \((r+(p))(z+p^jP)=rz+p^jP\) for \(z\in p^{j-1}P\).  This is
well-defined: changing \(z\) changes \(rz\) by an element of \(p^jP\), and
if \(r-r'=ps\), then \((r-r')z\in p^jP\).  Thus the quotient is an
\(R/(p)\)-module, hence a vector space.

\emph{Intrinsic character.}
If \(f:P\to Q\) is an isomorphism, then
\(f(p^kP)=p^kf(P)=p^kQ\) for every \(k\). Hence
\[
\overline f_j:p^{j-1}P/p^jP\longrightarrow p^{j-1}Q/p^jQ,
\qquad z+p^jP\longmapsto f(z)+p^jQ
\]
is well-defined and \(R/(p)\)-linear. The map induced in the same way by
\(f^{-1}\) is its inverse. Thus corresponding layers are isomorphic, and
\(a_j(P)\) depends only on \(P\).

\emph{Additivity.}
For modules \(P_1,P_2\), powers of \(p\) commute with their direct sum. There
is therefore a natural \(R/(p)\)-linear isomorphism
\[
\frac{p^{j-1}(P_1\oplus P_2)}{p^j(P_1\oplus P_2)}
\cong
\frac{p^{j-1}P_1}{p^jP_1}\oplus
\frac{p^{j-1}P_2}{p^jP_2}.
\]
Explicitly, it is given by
\[
(z_1,z_2)+p^j(P_1\oplus P_2)
\longmapsto (z_1+p^jP_1,z_2+p^jP_2)
\]
and its inverse is obtained by pairing representatives. Iterating this binary
isomorphism gives the corresponding statement for every finite direct sum;
Corollary~\ref{cor:finite-direct-sum-dimension} then gives
\[
a_j(P_1\oplus\cdots\oplus P_n)=\sum_{i=1}^n a_j(P_i).
\]

\emph{Contribution of one block.}
For \(R/(p^e)\), the \(j\)-th layer is zero when \(e<j\).  When \(e\geq j\),
it is generated by the class of \(p^{j-1}\), whose annihilator in that layer
is \((p)\); it is one copy of \(R/(p)\).  Additivity proves the formula.
\end{proof}

The decisive recovery formula is
\begin{equation}\label{eq:primary-layer-recovery}
m_j:=a_j(P)-a_{j+1}(P)
=\#\{\text{summands }R/(p^j)\text{ in }P\}.
\end{equation}
It converts intrinsic row lengths of the filtration into column heights,
hence into elementary divisors.

\begin{theorem}[Uniqueness: elementary-divisor form]
\label{thm:uniqueness-elementary-divisor}
For each prime representative \(p\) and each \(j\geq1\), the multiplicity of
\(R/(p^j)\) in a finitely generated PID module is uniquely determined by the
module. Hence the elementary-divisor decomposition is unique up to associates
and order; the free rank is unique as well.
\end{theorem}
\begin{proof}
The intrinsic module \(T(M)\) and
Theorem~\ref{thm:pid-structure-existence}(3) give
\(M/T(M)\cong R^r\), so IBN determines \(r\). Each primary component
\(M_p\) is intrinsic. On it, the intrinsic sequence \(a_j\)
determines through \eqref{eq:primary-layer-recovery} the multiplicity of every
block \(R/(p^j)\). Thus all elementary divisors are unique, prime by prime.
\end{proof}

Elementary divisors determine invariant factors by a canonical regrouping
procedure. Here is that regrouping explicitly. Suppose the elementary divisors are
\[
p,\quad p^3,\quad q^2,\quad q^4,\quad q^5,
\]
for nonassociate primes \(p,q\).  Arrange each prime's powers by increasing
exponent and pad the shorter list on the left by units:
\[
\begin{array}{ccc}
1&p&p^3\\
q^2&q^4&q^5
\end{array}.
\]
Multiplying down the columns gives
\[
q^2,\qquad pq^4,\qquad p^3q^5,
\]
and visibly \(q^2\mid pq^4\mid p^3q^5\).  Factoring an invariant-factor
chain reverses this construction. Thus elementary-divisor uniqueness already
provides one route to invariant-factor uniqueness; after establishing
cancellation, we will also give a direct proof in invariant-factor language.

\section{Cancellation and Consequences}
Cancellation is now a consequence of the intrinsic classification, rather
than an unproved step inside its proof.

\begin{corollary}[Cyclic prime-power cancellation]
Let \(M,N\) be finitely generated \(p\)-primary modules and let \(e\geq1\).
If
\[
M\oplus R/(p^e)\cong N\oplus R/(p^e),
\]
then \(M\cong N\).
\end{corollary}
\begin{proof}
Theorem~\ref{thm:uniqueness-elementary-divisor} says that the two sides have
the same multiplicity of every elementary divisor. Removing the one common
occurrence of \(p^e\) leaves identical elementary-divisor multisets for \(M\)
and \(N\).
\end{proof}

\begin{corollary}[Cyclic cancellation]\label{cor:cyclic-cancellation}
Let \(M,N\) be finitely generated torsion \(R\)-modules and let \(0\ne d\in R\).
If
\[
M\oplus R/(d)\cong N\oplus R/(d),
\]
then \(M\cong N\).
\end{corollary}
\begin{proof}
If \(d\) is a unit, the added cyclic module is zero and there is nothing to
prove. Otherwise factor \(d\), up to a unit, into powers of distinct fixed
prime representatives. The Chinese Remainder Theorem expresses \(R/(d)\) as
the direct sum of the corresponding prime-power cyclic modules.
Theorem~\ref{thm:uniqueness-elementary-divisor} says that the two sides of the
given isomorphism have the same multiset of elementary divisors. Remove the
common elementary divisors contributed by \(R/(d)\). The modules \(M\) and
\(N\) then have identical elementary-divisor decompositions, so \(M\cong N\).
\end{proof}

\begin{theorem}[Uniqueness: invariant-factor form]
\label{thm:uniqueness-invariant-factor}
Suppose a finitely generated torsion \(R\)-module has two decompositions
\[
M\cong R/(a_1)\oplus\cdots\oplus R/(a_r)
\cong R/(b_1)\oplus\cdots\oplus R/(b_s),
\]
where
\[
a_1\mid\cdots\mid a_r,
\qquad b_1\mid\cdots\mid b_s,
\]
and every displayed cyclic factor is nonzero. Then \(r=s\) and, after
matching the two divisibility chains, \((a_i)=(b_i)\) for every \(i\).
\end{theorem}
\begin{proofidea}
The annihilator identifies the largest invariant factor. Cyclic cancellation
removes that factor, and induction repeats the argument on the remaining
chains.
\end{proofidea}
\begin{proof}
We induct on \(r+s\). If \(r=0\), then \(M=0\); because every displayed
factor in the second decomposition is nonzero, necessarily \(s=0\). The same
argument applies with the two decompositions reversed, establishing the base
case.

Now suppose \(r,s>0\). The annihilator of a finite direct sum is the
intersection of the annihilators of its summands, because a scalar
annihilates the direct sum exactly when it annihilates every coordinate. The
divisibility chains give
\[
\operatorname{Ann}(M)
=(a_1)\cap\cdots\cap(a_r)=(a_r),
\qquad
\operatorname{Ann}(M)
=(b_1)\cap\cdots\cap(b_s)=(b_s).
\]
Thus \((a_r)=(b_s)\), so \(R/(a_r)\cong R/(b_s)\). Put
\[
A=\bigoplus_{i<r}R/(a_i),
\qquad
B=\bigoplus_{j<s}R/(b_j).
\]
After identifying the last two factors with the common cyclic type
\(R/(a_r)\), the original isomorphisms give
\[
A\oplus R/(a_r)\cong B\oplus R/(a_r).
\]
Corollary~\ref{cor:cyclic-cancellation} yields \(A\cong B\). The total number
of displayed factors has decreased, so induction gives \(r-1=s-1\) and
\((a_i)=(b_i)\) for \(i<r\). Together with \((a_r)=(b_s)\), this proves
\(r=s\) and equality of all matched ideals.
\end{proof}

The two uniqueness theorems reflect the two normal forms. Primary layers
recover all prime-power block multiplicities simultaneously and prove
Theorem~\ref{thm:uniqueness-elementary-divisor}. In contrast, the direct
invariant-factor argument follows the chain
\[
\text{annihilator}\longrightarrow\text{largest invariant factor}
\longrightarrow\text{cancellation}\longrightarrow\text{induction},
\]
detecting and removing one factor at a time. Consequently, invariant factors
are unique up to associates, while elementary divisors are unique up to
associates and order.

\begin{remark}
The annihilator alone does not single out a preferred embedded cyclic summand
inside an abstract direct sum, so cancellation cannot be assumed. The
primary-layer argument is the natural intrinsic proof of uniqueness in
elementary-divisor form. Once cyclic cancellation has been proved from that
classification, the annihilator argument gives a particularly short inductive
proof of uniqueness in invariant-factor form.
\end{remark}

\section{Finitely Generated Abelian Groups}
The familiar classification of finitely generated abelian groups is now a
corollary, rather than a separate theorem: abelian groups are precisely
\(\mathbb Z\)-modules.

\begin{corollary}[Finitely generated abelian groups]
Every finitely generated abelian group is uniquely isomorphic to
\[
\mathbb Z^r\oplus\mathbb Z/(d_1)\oplus\cdots\oplus\mathbb Z/(d_k),
\qquad 1<d_1\mid\cdots\mid d_k,
\]
or, equivalently, to a direct sum of a free abelian group and cyclic
prime-power groups.
\end{corollary}

Indeed, abelian groups are precisely \(\mathbb Z\)-modules and
\(\mathbb Z\) is a PID, so this is the PID module structure theorem applied
to \(R=\mathbb Z\).  Factoring the invariant factors into prime powers gives
the elementary-divisor form.

\begin{corollary}[Finiteness and order]
A finitely generated abelian group is finite if and only if its free rank is
zero.  In that case,
\[
\left|\mathbb Z/(d_1)\oplus\cdots\oplus\mathbb Z/(d_k)\right|
=d_1\cdots d_k.
\]
\end{corollary}
\begin{proof}
A nonzero free summand \(\mathbb Z\) makes the group infinite.  With no free
summand, the group is a finite Cartesian product of finite cyclic groups, so
its order is the product of their orders.
\end{proof}

\begin{example}[Free and torsion parts]
The group
\[
G=\mathbb Z^2\oplus\mathbb Z/6\mathbb Z\oplus\mathbb Z/30\mathbb Z
\]
has free rank \(2\), torsion subgroup of order \(180\), and is infinite.
Its elementary divisors are
\[
2,\ 3,\ 2,\ 3,\ 5
\]
with prime powers grouped by their two invariant-factor columns; more
transparently,
\[
G_{\mathrm{tors}}\cong
\mathbb Z/2\oplus\mathbb Z/3\oplus
\mathbb Z/2\oplus\mathbb Z/3\oplus\mathbb Z/5.
\]
\end{example}

\begin{example}[Converting the two forms]
\(\mathbb Z/12\mathbb Z\oplus\mathbb Z/18\mathbb Z\) has
elementary divisors \(\mathbb Z/4\mathbb Z\), \(\mathbb Z/3\mathbb Z\),
\(\mathbb Z/2\mathbb Z\), and \(\mathbb Z/9\mathbb Z\); regrouping gives
the invariant factors
\[
\mathbb Z/6\mathbb Z\oplus\mathbb Z/36\mathbb Z.
\]
Its order is \(6\cdot36=216\), the same as \(12\cdot18\).
\end{example}

\begin{example}[Classification at a fixed order]
Every abelian group of order \(72=2^3 3^2\) is obtained by independently
choosing a partition of \(3\) for the \(2\)-primary part and a partition of
\(2\) for the \(3\)-primary part.  The three possibilities
\[
\mathbb Z/8,\qquad \mathbb Z/4\oplus\mathbb Z/2,\qquad
(\mathbb Z/2)^3
\]
for the \(2\)-part combine with the two possibilities
\[
\mathbb Z/9,\qquad(\mathbb Z/3)^2
\]
for the \(3\)-part.  Thus there are exactly six isomorphism types, one for
each pairing of the displayed choices.  For example, pairing the middle and
second choices gives
\[
\mathbb Z/4\oplus\mathbb Z/2\oplus
\mathbb Z/3\oplus\mathbb Z/3.
\]
No search through group operations is needed.
\end{example}

\section{Smith Normal Form}
Theorem~\ref{thm:pid-submodules-compatible-bases}(2) gives the conceptual
existence argument, while the preceding uniqueness theorems identify the
elementary-divisor and invariant-factor data intrinsically. Smith normal form
is the matrix manifestation of the same structure. We first explain why it exists, then
show how B\'ezout row and column operations compute it.

For matrices \(A,B\in M_{n\times m}(R)\), say that \(A\) and \(B\) are
\emph{equivalent} if \(B=UAV\) for some
\(U\in\operatorname{GL}_n(R)\) and \(V\in\operatorname{GL}_m(R)\).
Here \(U\) changes the basis of the target \(R^n\), while \(V\) changes the
basis of the source \(R^m\). This differs fundamentally from similarity,
\(B=P^{-1}AP\), which concerns square matrices representing one endomorphism
after one change of basis. Thus matrix equivalence is not operator similarity.
This distinction matters when Smith reduction is applied to \(xI-A\), whereas
rational canonical form classifies \(A\) up to similarity.

\begin{definition}
A matrix over a PID is in \emph{Smith normal form} if it is
\[\operatorname{diag}(d_1,\ldots,d_t,0,\ldots,0),\qquad
d_1\mid\cdots\mid d_t,\quad d_i\ne0.\]
\end{definition}

The elementary invertible row and column operations are: interchange two
rows or columns; multiply one row or column by a unit; and add a scalar
multiple of one row or column to another. Each is multiplication on the
appropriate side by an elementary matrix. Its inverse is respectively the
same interchange, multiplication by the inverse unit, or addition of the
negative scalar multiple, so each operation preserves matrix equivalence.

\begin{theorem}[Smith normal form]
For every matrix \(A\) over a PID, there are invertible matrices \(U,V\)
such that \(UAV\) is in Smith normal form.
\end{theorem}
\begin{proofidea}
Regard the matrix as a homomorphism \(A:R^m\to R^n\). Choose compatible
target coordinates for \(\operatorname{im}(A)\), then split the surjection from the source
onto that free image.  Compatible source and target bases make the matrix
diagonal.
\end{proofidea}
\begin{proof}
Apply Theorem~\ref{thm:pid-submodules-compatible-bases}(2) to
\(\operatorname{im}(A)\subseteq R^n\).  There is a target basis
\(y_1,\ldots,y_n\) such that
\[
\operatorname{im}(A)=Rd_1y_1\oplus\cdots\oplus Rd_ty_t,
\qquad d_1\mid\cdots\mid d_t.
\]
The displayed generators form a basis of \(\operatorname{im}(A)\).  Because
this module is free, it is projective by Chapter~5, so the surjection
\[
R^m\twoheadrightarrow\operatorname{im}(A)
\]
splits.  Choose a section \(s\), put \(x_i=s(d_iy_i)\), and choose a basis
of \(\ker A\), which is free by Section~\ref{sec:pid-submodules-free}.  The
splitting gives
\[
R^m=s(\operatorname{im}(A))\oplus\ker A,
\]
so \(x_1,\ldots,x_t\), followed by the kernel basis, is a basis of the
source.  In these bases,
\[
A(x_i)=d_iy_i\quad(1\leq i\leq t),
\qquad A(\ker A)=0,
\]
and the matrix is \(\operatorname{diag}(d_1,\ldots,d_t,0,\ldots,0)\).
The two changes of basis give the required invertible \(U,V\).
\end{proof}

\subsection*{Constructive viewpoint: B\'ezout row and column reduction}
The structural proof explains why Smith form exists.  The next argument
explains how to reach it through explicit invertible operations.

\begin{proposition}[Constructive PID diagonal reduction]
Every finite matrix over a PID can be transformed to Smith normal form by
invertible row and column operations.
\end{proposition}
\begin{proofidea}
Improve the upper-left pivot until it divides every entry; prove termination
by ACC on principal ideals; clear its row and column; then induct on the
remaining submatrix.
\end{proofidea}
\begin{proof}
If the matrix is zero, there is nothing to prove.  Otherwise move a nonzero
entry \(d\) to the upper-left corner.  Suppose first that an entry \(b\) in
the first column is not divisible by \(d\).  Let \(g\) generate \((d,b)\),
write \(d=gd'\), \(b=gb'\), and choose \(u,v\in R\) with \(ud+vb=g\).  Then
\[
\begin{pmatrix}u&v\\-b'&d'\end{pmatrix}
\begin{pmatrix}d\\b\end{pmatrix}
=\begin{pmatrix}g\\0\end{pmatrix},
\qquad
\det\begin{pmatrix}u&v\\-b'&d'\end{pmatrix}=ud'+vb'=1.
\]
This invertible row operation replaces the pivot by \(g\).  Analogous right
multiplication on two columns handles a nondivisible first-row entry.

Repeat these B\'ezout improvements until the pivot divides every entry in its
row and column, then clear the rest of that row and column.  The matrix has
the form
\[
\begin{pmatrix}d&0\\0&A_1\end{pmatrix}.
\]
If an interior entry \(b\) of \(A_1\) is not divisible by \(d\), add its row
to the first row.  The first column remains unchanged, while \(b\) enters the
first row.  Apply to the pivot column and the column containing \(b\) the
two-column operation
\[
\begin{pmatrix}u&-b'\\v&d'\end{pmatrix},
\qquad ud+vb=g,\quad d=gd',\quad b=gb'.
\]
Indeed,
\[
(d,b)\begin{pmatrix}u&-b'\\v&d'\end{pmatrix}=(g,0),
\]
and its determinant is \(1\).  The new pivot generates \((d,b)\), and
\(d\nmid b\) implies
\[
(d)\subsetneq(d,b).
\]
Every such improvement strictly enlarges the principal pivot ideal.
Chapter~4 proves that a PID is Noetherian and that ACC terminates ascending
ideal chains, so this process stops.

At termination, clear the first row and column once more.  No interior entry
can fail to be divisible by the pivot, or the preceding step would improve it
again.  Thus the first diagonal entry \(d_1\) divides every entry of the
remaining matrix.  Apply induction to that submatrix.  Its operations form
only linear combinations of its entries, so divisibility by \(d_1\) survives,
and the resulting diagonal satisfies
\[
d_1\mid d_2\mid\cdots\mid d_t.
\]
Every operation used is invertible; no Euclidean division is required.
\end{proof}

Over an arbitrary PID this is an explicit terminating reduction procedure,
with termination supplied by ACC. Effectiveness additionally requires the
ability to compute gcds and B\'ezout coefficients. Over Euclidean domains such
as \(\mathbb Z\) and \(F[x]\), the Euclidean algorithm supplies these, so the
procedure is a computational algorithm in the usual sense.

Presentations explain what this computation does to a module.  Chosen
generators and relations give an exact sequence
\[
\begin{tikzcd}[column sep=large]
R^m\arrow[r,"A"]&R^n\arrow[r,two heads]&M\arrow[r]&0,
\end{tikzcd}
\]
where the columns of \(A\) record the relations and
\(M\cong\operatorname{coker}(A)\).  Row operations change the target basis,
column operations change the source basis, and neither changes the cokernel.
Smith reduction therefore separates the defining relations coordinatewise.

\begin{corollary}[Kernel and cokernel from Smith form]
\label{cor:smith-kernel-cokernel}
Suppose \(A:R^m\to R^n\) has Smith form
\(\operatorname{diag}(d_1,\ldots,d_t,0,\ldots,0)\). Then
\[
\operatorname{rank}(\operatorname{im}A)=t,
\qquad \ker A\cong R^{m-t},
\]
and
\[
\operatorname{coker}A\cong
R/(d_1)\oplus\cdots\oplus R/(d_t)\oplus R^{n-t}.
\]
Unit \(d_i\) contribute zero cyclic summands and may be omitted. Thus Smith
reduction diagonalizes the relations in a finite presentation.
\end{corollary}
\begin{proof}
Equivalent matrices represent the same map after source and target
isomorphisms. For the diagonal map the assertions follow coordinatewise;
multiplication by each nonzero \(d_i\) is injective because a PID is a domain.
\end{proof}

Over a field every nonzero Smith entry is a unit and can be normalized to
\(1\). Thus Smith form becomes
\(\operatorname{diag}(1,\ldots,1,0,\ldots,0)\): ordinary rank normal form is
the field case of the divisibility data carried by Smith form over a PID.

If \(A\) is \(n\times n\) and has full rank, determinant multiplicativity
gives \(\det A\sim d_1\cdots d_n\). For an integral matrix with
\(\det A\ne0\), the cokernel formula therefore gives
\[
\lvert\mathbb Z^n/A\mathbb Z^n\rvert=\lvert\det A\rvert.
\]

\noindent\textbf{Worked example.} Over \(\mathbb Z\), elementary operations
give
\[
\begin{pmatrix}2&4\\6&8\end{pmatrix}
\xrightarrow{R_2\mapsto R_2-3R_1}
\begin{pmatrix}2&4\\0&-4\end{pmatrix}
\xrightarrow{C_2\mapsto C_2-2C_1}
\begin{pmatrix}2&0\\0&-4\end{pmatrix}.
\]
Multiplying the second row by \(-1\) yields
\[
\operatorname{diag}(2,4).
\]
Consequently its cokernel is
\[
\mathbb Z/2\mathbb Z\oplus\mathbb Z/4\mathbb Z.
\]
Equivalently, the abelian group with two generators and relation columns
\((2,6)^t\) and \((4,8)^t\) has precisely these invariant factors.  This is
the promised passage from a presentation matrix to the classification of its
cokernel.

\noindent\textbf{A genuine B\'ezout pivot improvement.}
For the column matrix \(A=(6,15)^t\), the pivot \(6\) does not divide \(15\).
Since \((-2)6+15=3\),
\[
\begin{pmatrix}-2&1\\-5&2\end{pmatrix}
\begin{pmatrix}6\\15\end{pmatrix}
=\begin{pmatrix}3\\0\end{pmatrix},
\qquad
\det\begin{pmatrix}-2&1\\-5&2\end{pmatrix}=1.
\]
This invertible row operation replaces the pivot by
\(\gcd(6,15)=3\), displaying the B\'ezout improvement that repeated
subtraction can conceal.

For a matrix, let \(D_j(A)\) be the ideal generated by its \(j\times j\)
minors, with \(D_j(A)=R\) for \(j\leq0\) and \(D_j(A)=0\) when \(j\)
exceeds the possible minor size. Elementary row and column operations preserve
each \(D_j(A)\), because each new minor is an \(R\)-linear combination of old
minors, and the inverse operation gives the reverse inclusion.  In Smith form
with \(t\) nonzero entries,
\[
D_j(A)=(d_1\cdots d_j)\quad(1\leq j\leq t),
\qquad D_j(A)=0\quad(j>t).
\]
If \(D_j(A)=(\Delta_j)\), then, up to associates,
\(d_1\sim\Delta_1\) and
\(\Delta_j\sim\Delta_{j-1}d_j\) successively. Concretely, \(d_1\) is
associate to the gcd of all entries, \(d_1d_2\) to the gcd of all
\(2\times2\) minors, and so forth.

\begin{corollary}[Uniqueness of Smith normal form]
\label{cor:smith-uniqueness}
Suppose
\[
\operatorname{diag}(d_1,\ldots,d_t,0,\ldots,0),\qquad
\operatorname{diag}(e_1,\ldots,e_s,0,\ldots,0)
\]
are Smith forms equivalent to the same matrix, with both nonzero lists ordered
by divisibility. Then \(t=s\) and \(d_i\sim e_i\) for every \(i\).
With fixed associate representatives, such as positive integers in
\(\mathbb Z\) or monic polynomials in \(F[x]\), the Smith form is literally
unique.
\end{corollary}
\begin{proof}
Equivalent matrices have the same determinantal ideals. The integer \(t\) is
the largest \(j\geq0\) for which \(D_j(A)\ne0\), with \(D_0(A)=R\) by
convention. Hence equivalent matrices have the same \(t\). Then
\((d_1\cdots d_j)=(e_1\cdots e_j)\) for \(1\leq j\leq t\); cancellation in the
domain, applied successively, gives \((d_j)=(e_j)\).
\end{proof}

\subsection*{Alternative uniqueness via Smith normal form and Fitting ideals}
Smith uniqueness compares equivalent matrices. An abstract module, however,
can have many finite presentations, so another step is required before minors
become module invariants.

Let
\[
R^m\xrightarrow{A}R^n\longrightarrow M\longrightarrow0
\]
be a finite presentation. For every integer \(k\), define
\[
\operatorname{Fitt}_k(M)=D_{n-k}(A),
\]
using \(D_j(A)=R\) for \(j\leq0\) and \(D_j(A)=0\) beyond the possible
minor size.

\begin{theorem}[Presentation independence of Fitting ideals]
\label{thm:fitting-presentation-independent}
The ideal \(\operatorname{Fitt}_k(M)\) is independent of the chosen finite
presentation. Hence every Fitting ideal is intrinsic to \(M\).
\end{theorem}
\begin{proof}
We compare separately the relations and the module generators.

First fix a surjection \(R^n\twoheadrightarrow M\). If the columns of two
matrices \(A\) and \(A'\) generate its kernel, each column of either matrix is
a linear combination of columns of the other. Thus \(A'=AC\) and \(A=A'C'\)
for suitable matrices. Fix a \(j\times j\) minor of \(AC\). After selecting
its \(j\) rows, each of its columns is an \(R\)-linear combination of the
columns of \(A\) restricted to those rows. Multilinearity expands its
determinant as an \(R\)-linear combination of determinants formed from \(j\)
restricted columns of \(A\). Terms with a repeated column vanish; every
remaining determinant is, up to sign, a \(j\times j\) minor of \(A\).
Therefore \(D_j(AC)\subseteq D_j(A)\), so
\(D_j(A')\subseteq D_j(A)\). The reverse factorization gives equality, and
the ideals do not depend on the chosen finite generating list of relations.

Now adjoin to module generators \(y_1,\ldots,y_n\) one redundant generator
\(y=\sum c_i y_i\). If \(c=(c_1,\ldots,c_n)^t\), a relation matrix for the
enlarged list is
\[
B=\begin{pmatrix}A&-c\\0&1\end{pmatrix}.
\]
Indeed, \((u,s)\) is a relation on the enlarged list exactly when
\(u+sc\in\operatorname{im}A\), which is exactly when
\((u,s)=B(z,s)\) for some \(z\).
Adding \(c_i\) times the last row to row \(i\) makes \(B\) equivalent to
\(A\oplus(1)\). A Laplace expansion of any \(j\times j\) minor of \(A\)
along one row expresses it as an \(R\)-linear combination of
\((j-1)\times(j-1)\) minors of \(A\). Hence
\(D_j(A)\subseteq D_{j-1}(A)\). A \(j\)-minor of the block matrix either avoids the unit
block or uses it, so
\[
D_j(A\oplus(1))=D_j(A)+D_{j-1}(A)=D_{j-1}(A),
\]
because \(D_j(A)\subseteq D_{j-1}(A)\). Therefore
\[
D_{(n+1)-k}(B)=D_{n-k}(A),
\]
so adjoining a redundant module generator does not change
\(\operatorname{Fitt}_k\).

Finally, take two finite generating lists for \(M\) and adjoin every member
of each list to the other. Both become the same combined list, one redundant
generator at a time. For that common surjection, the first paragraph compares
any two finite relation lists. Such lists exist because its kernel is a
submodule of a finite free module over the PID \(R\), hence is finitely
generated by Theorem~\ref{thm:pid-submodules-compatible-bases}. Thus the two
original presentations give the same ideals.
\end{proof}

\begin{theorem}[Invariant factors from Fitting ideals]
\label{thm:fitting-invariant-factors}
Suppose existence gives
\[
M\cong R^r\oplus R/(d_1)\oplus\cdots\oplus R/(d_t),
\qquad d_1\mid\cdots\mid d_t,
\]
where the \(d_i\) are nonzero nonunits. Then the free rank and all associate
classes \((d_i)\) are determined by \(M\).
\end{theorem}
\begin{proof}
The free rank \(r\) is intrinsic because
\(M/T(M)\cong R^r\) and IBN determines \(r\). A diagonal presentation with
\(r+t\) generators has relation matrix consisting of a zero \(r\)-row block
above \(\operatorname{diag}(d_1,\ldots,d_t)\). Consequently
\[
\begin{aligned}
\operatorname{Fitt}_k(M)&=0 &&(k<r),\\
\operatorname{Fitt}_{r+t-j}(M)&=(d_1\cdots d_j)
&&(1\leq j\leq t),\\
\operatorname{Fitt}_{r+t}(M)&=R.
\end{aligned}
\]
Because every \(d_i\) is a nonunit, \(r+t\) is the least index at which the
Fitting ideal becomes \(R\); hence \(t\) is intrinsic. Reading the middle
line from \(j=1\) upward gives
\[
(d_1),\quad(d_1d_2),\quad\ldots,\quad(d_1\cdots d_t).
\]
Successive cancellation recovers each \((d_j)\). This proves invariant-factor
uniqueness without choosing or identifying a summand of \(M\). Factoring the
recovered invariant factors into prime powers then recovers the elementary
divisors as well.
\end{proof}

The three uniqueness arguments illuminate different structure. Primary
layers use the intrinsic \(p\)-primary filtration and recover all prime-power
block multiplicities at once; this remains the conceptual main proof.
Annihilator plus cancellation detects the largest invariant factor, removes
it, and inducts. Smith form and Fitting ideals pass from a finite presentation
to minors, prove those ideals presentation-independent, and recover invariant
factors. These are complementary rather than circular: the Fitting argument
uses only existence of a decomposition, not its uniqueness, and the Smith
uniqueness corollary rests directly on determinantal ideals. Fitting ideals
also make sense beyond PIDs, where they help measure how many generators a
module needs locally.

\section{Linear Operators as \texorpdfstring{\(F[x]\)}{F[x]}-Modules}
\label{sec:linear-operators-fx-modules}
Chapter~5 showed what generalized eigenvectors and Jordan chains look like.
It did not yet explain why enough compatible chains exist or why their block
sizes are intrinsic. Why do the canonical blocks exist, and why are their
sizes unique? The structural answer is to let \(x\) act as \(T\). Then \(V\)
becomes a finitely generated torsion \(F[x]\)-module, and the PID structure
theorem answers both questions: the blocks are its cyclic summands, whose
types are intrinsic.

Let \(T:V\to V\) be linear over a field \(F\). Define
\[p(x)\cdot v=p(T)v,\qquad
p(T)=a_0I+a_1T+\cdots+a_nT^n.\]
Because \((pq)(T)=p(T)q(T)\), the module axioms hold.
Write \(V_T\) when the operator defining this \(F[x]\)-module structure
needs to be displayed.

Recall from Section~\ref{sec:matrices-change-basis} the direct sums of
operators and their block-diagonal matrices. Direct calculation on pairs
gives, for every \(p\in F[x]\),
\[
p(T_1\oplus T_2)=p(T_1)\oplus p(T_2),
\]
and the identity map on the underlying direct sum is an \(F[x]\)-module
isomorphism
\[
(V_1\oplus V_2)_{T_1\oplus T_2}
\cong (V_1)_{T_1}\oplus(V_2)_{T_2}.
\]

\begin{proposition}
A subspace \(W\subseteq V\) is \(T\)-invariant if and only if it is an
\(F[x]\)-submodule for this action.
\end{proposition}
\begin{proof}
If \(W\) is a submodule, then \(Tw=x\cdot w\in W\). Conversely, if
\(T(W)\subseteq W\), every polynomial in \(T\) preserves \(W\), so \(W\)
is an \(F[x]\)-submodule.
\end{proof}

\begin{proposition}[Intertwiners are module maps]\label{prop:intertwiner-module-map}
Let \(T:V\to V\) and \(S:W\to W\).  An \(F\)-linear map
\(\phi:V\to W\) is \(F[x]\)-linear from \(V_T\) to \(W_S\) if and only if
\[
\phi T=S\phi.
\]
Consequently \(V_T\cong W_S\) as \(F[x]\)-modules if and only if there is
an invertible intertwining map between \(T\) and \(S\).  In coordinates this
is precisely similarity of their matrices.
\end{proposition}
\begin{proof}
If \(\phi\) is \(F[x]\)-linear, then
\[
\phi(Tv)=\phi(x\cdot v)=x\cdot\phi(v)=S\phi(v).
\]
Conversely, \(\phi T=S\phi\) implies \(\phi T^j=S^j\phi\) by induction,
and linear combinations give
\(\phi p(T)=p(S)\phi\) for every polynomial \(p\).  Thus \(\phi\) is
\(F[x]\)-linear.  If \(\phi\) is invertible and its matrix is \(P\), the
intertwining equation becomes \(PA=BP\), or \(B=PAP^{-1}\), exactly the
change-of-basis relation.
\end{proof}

\begin{proposition}[Finite-dimensional operator modules are torsion]
\label{prop:operator-module-torsion}
If \(V\) is finite-dimensional, then \(V_T\) is a finitely generated torsion
\(F[x]\)-module.
\end{proposition}
\begin{proof}
An \(F\)-basis \(v_1,\ldots,v_n\) generates \(V_T\), because constants in
\(F[x]\) already produce every \(F\)-linear combination.  For each \(v_i\),
the \(n+1\) vectors
\[
v_i,Tv_i,\ldots,T^nv_i
\]
are linearly dependent, so a nonzero polynomial \(f_i\) satisfies
\(f_i(T)v_i=0\).  The nonzero product \(f_1\cdots f_n\) annihilates every
basis vector and hence all of \(V\).  In particular every vector is torsion.
\end{proof}

For a monic polynomial
\[
d(x)=x^m+a_{m-1}x^{m-1}+\cdots+a_0
\]
define its \emph{companion matrix} by
\[
C_d=
\begin{pmatrix}
0&0&\cdots&0&-a_0\\
1&0&\cdots&0&-a_1\\
0&1&\cdots&0&-a_2\\
\vdots&&\ddots&\vdots&\vdots\\
0&0&\cdots&1&-a_{m-1}
\end{pmatrix}.
\]
The next proposition derives this matrix directly from the first polynomial
relation in an operator orbit.

\begin{proposition}[Cyclic subspaces, companion matrices, and annihilators]
\label{prop:cyclic-subspace-annihilator}
Let \(0\ne v\in V\), and let
\[
W=F[x]v=\operatorname{span}_F\{v,Tv,T^2v,\ldots\}.
\]
If \(\dim W=m\), then
\[
\beta=(v,Tv,\ldots,T^{m-1}v)
\]
is a basis of \(W\). There is a unique monic polynomial
\[
d(x)=x^m+a_{m-1}x^{m-1}+\cdots+a_1x+a_0
\]
generating \(\operatorname{Ann}_{F[x]}(v)\), and
\[
F[x]v\cong F[x]/(d),
\qquad
[T|_W]_\beta=C_d.
\]
\end{proposition}
\begin{proof}
Let \(k\) be the least index for which
\(v,Tv,\ldots,T^kv\) are linearly dependent. Then the first \(k\) vectors
are independent, and the first dependence expresses \(T^kv\) in their span.
Applying \(T\) repeatedly shows that every later orbit vector remains in that
span. Thus those \(k\) vectors form a basis of \(W\), so \(k=\dim W=m\).
Equivalently, \(v,Tv,\ldots,T^{m-1}v\) are independent, and adjoining
\(T^mv\) produces the first dependence.

Consequently there are scalars \(c_0,\ldots,c_m\in F\) such that
\[
c_mT^mv+c_{m-1}T^{m-1}v+\cdots+c_1Tv+c_0v=0.
\]
Here \(c_m\ne0\): otherwise this would already be a dependence among the
first \(m\) orbit vectors. Because \(F\) is a field, \(c_m\) is invertible.
We intentionally divide the entire relation by \(c_m\), normalizing its
leading coefficient to \(1\). With \(a_i=c_i/c_m\), the relation becomes
\[
T^mv+a_{m-1}T^{m-1}v+\cdots+a_1Tv+a_0v=0,
\]
or equivalently
\[
T^mv=-a_0v-a_1Tv-\cdots-a_{m-1}T^{m-1}v.
\]
This deliberate normalization makes
\(d(x)=x^m+a_{m-1}x^{m-1}+\cdots+a_1x+a_0\) monic, matching the convention
for generators of ideals in \(F[x]\).

On the orbit basis, the action is
\[
T(v)=Tv,\quad T(Tv)=T^2v,\quad\ldots,\quad
T(T^{m-2}v)=T^{m-1}v,
\]
while the last basis vector satisfies
\[
T(T^{m-1}v)=-a_0v-a_1Tv-\cdots-a_{m-1}T^{m-1}v.
\]
Reading these coordinate columns gives
\[
[T|_W]_\beta=C_d=
\begin{pmatrix}
0&0&\cdots&0&-a_0\\
1&0&\cdots&0&-a_1\\
0&1&\cdots&0&-a_2\\
\vdots&&\ddots&\vdots&\vdots\\
0&0&\cdots&1&-a_{m-1}
\end{pmatrix}.
\]
Thus a companion matrix is precisely the matrix of \(T\) on a cyclic
\(T\)-invariant subspace in its natural orbit basis.

The normalized relation says \(d(T)v=0\), so
\(d\in\operatorname{Ann}_{F[x]}(v)\). No nonzero polynomial of degree less
than \(m\) can annihilate \(v\), since it would give a dependence among
\(v,Tv,\ldots,T^{m-1}v\). If \(d_v\) is the monic generator of
\(\operatorname{Ann}(v)\), then \(d_v\mid d\), while the preceding
independence gives \(\deg d_v\geq m\). Hence \(d_v=d\) and
\(\operatorname{Ann}(v)=(d)\).

Finally, the surjective module map
\[
F[x]\longrightarrow F[x]v,\qquad p\longmapsto p(T)v,
\]
has kernel \((d)\). The First Isomorphism Theorem gives
\(F[x]v\cong F[x]/(d)\), making the concrete orbit calculation and the
abstract cyclic-module description identical.
\end{proof}

\begin{theorem}[The \(xI-A\) presentation]
\label{thm:operator-cokernel-presentation}
Let \(V\) have ordered basis \(\beta=(v_1,\ldots,v_n)\), let
\(T\in\mathcal L(V)\), and put \(A=[T]_\beta\). Then, as \(F[x]\)-modules,
\[
V_T\cong\operatorname{coker}\bigl(xI_n-A:F[x]^n\to F[x]^n\bigr).
\]
If the Smith form over \(F[x]\) is
\[
\operatorname{diag}(1,\ldots,1,d_1,\ldots,d_r),
\qquad d_1\mid\cdots\mid d_r,
\]
with the \(d_i\) monic nonunits, then they are precisely the invariant factors
of \(V_T\).
\end{theorem}
\begin{proof}
Send the standard basis vector \(e_i\) of \(F[x]^n\) to \(v_i\). This gives
a surjective \(F[x]\)-linear map to \(V_T\). If \(A=(a_{ij})\), column \(j\)
of \(xI_n-A\) maps to
\[
xv_j-\sum_i a_{ij}v_i=Tv_j-Tv_j=0,
\]
so the map factors through the cokernel. Conversely, send
\(\sum_i c_iv_i\) to the class of the constant vector \((c_i)\). In the
cokernel, the displayed column relations say
\(x[e_j]=\sum_i a_{ij}[e_i]\), so this inverse candidate respects the action
of \(x\), hence every polynomial. The two induced maps are inverse on the
displayed generators.

Smith equivalence preserves the cokernel. Corollary~\ref{cor:smith-kernel-cokernel}
therefore identifies it with the direct sum of the \(F[x]/(d_i)\); unit
entries contribute zero modules. Also
\(\det(xI_n-A)=\chi_T(x)\) is a nonzero monic polynomial, so \(xI_n-A\)
has full rank and its Smith form has no zero diagonal entries. The remaining
nonunit entries are exactly the invariant factors.
\end{proof}

Thus
\[
\operatorname{SNF}_{F[x]}(xI-A)
\longrightarrow\text{invariant factors of }V_T
\longrightarrow\text{rational canonical form}.
\]
The matrix \(xI-A\) is a matrix over \(F[x]\), and its Smith equivalence is
the computational bridge from presentations to similarity classification.

The cyclic-subspace proposition is the first half of the canonical-form dictionary:
a cyclic module \(F[x]/(d)\) is a companion block.  The special case
\(d=(x-\lambda)^k\) is the Jordan-chain block described below.

\[
\begin{array}{c|c}
\text{Linear algebra}&F[x]\text{-module theory}\\ \hline
T\text{-invariant subspace}&F[x]\text{-submodule}\\
\text{cyclic \(T\)-subspace generated by }v&F[x]v\\
p(T)v=0&p(x)\text{ annihilates }v\\
\text{similarity of operators}&F[x]\text{-module isomorphism}\\
\text{generalized eigenspace for }\lambda&(x-\lambda)\text{-primary component}\\
\text{companion block }C_d&F[x]/(d)\\
\text{Jordan block }J_k(\lambda)&F[x]/((x-\lambda)^k)
\end{array}
\]
This dictionary synthesizes the correspondences proved in this section and
turns the linear-algebra questions into the PID classification problem.

\begin{definition}
For an irreducible polynomial \(p\in F[x]\), the \emph{\(p\)-primary
component} of \(V_T\) is
\[
V_p=\{v\in V:p(T)^Nv=0\text{ for some }N\geq1\}.
\]
We also write this component as \((V_T)_p\) when the ambient operator module
needs to be emphasized.
\end{definition}

\begin{proposition}[Primary components and generalized eigenspaces]
\label{prop:primary-generalized-eigenspace}
For \(p(x)=x-\lambda\),
\[
V_p=G_\lambda(T).
\]
Thus the concrete generalized eigenspace of Chapter~5 is exactly the
\((x-\lambda)\)-primary component of the operator module.
\end{proposition}
\begin{proof}
The defining condition \(p(T)^Nv=0\) is precisely
\((T-\lambda I)^Nv=0\).  Taking all such vectors gives the same set in the
two definitions.
\end{proof}

If \(p_1,\ldots,p_s\) are the distinct monic irreducibles occurring in the
elementary divisors of \(V_T\), the PID primary decomposition gives
\[
V_T=(V_T)_{p_1}\oplus\cdots\oplus(V_T)_{p_s}.
\]

\begin{proposition}[Jordan-chain dictionary]
\label{prop:jordan-chain-cyclic-module}
A Jordan chain \(v_1,\ldots,v_k\) for \(\lambda\) spans a cyclic
\(F[x]\)-module isomorphic to
\[
F[x]/((x-\lambda)^k).
\]
Conversely, every such cyclic summand supplies a Jordan chain of length
\(k\).
\end{proposition}
\begin{proof}
The last vector \(v_k\) generates the chain because
\[
(x-\lambda)^{k-j}\cdot v_k=v_j.
\]
It is killed by \((x-\lambda)^k\), but not by a smaller positive power,
so if \(\operatorname{Ann}(v_k)=(d)\) with \(d\) monic, then
\(d\mid(x-\lambda)^k\). Hence \(d=(x-\lambda)^j\) for some \(j\leq k\).
No smaller positive power kills \(v_k\), so \(j=k\) and
\(\operatorname{Ann}(v_k)=((x-\lambda)^k)\). Conversely, in
\(F[x]/((x-\lambda)^k)\), use the ordered basis
\[
(x-\lambda)^{k-1},(x-\lambda)^{k-2},\ldots,x-\lambda,1.
\]
Multiplication by \(x-\lambda\) sends each basis vector to the preceding
one and kills the first.  Multiplication by \(x\) therefore has matrix
\(J_k(\lambda)\), and the basis is a Jordan chain in the convention of
Chapter~5.
\end{proof}

\section{Minimal Polynomial and Cayley--Hamilton}
The characteristic polynomial introduced in Chapter~5 detects eigenvalues.
The minimal polynomial records the smallest polynomial relation satisfied by
the whole operator and will control the sizes of canonical-form blocks.  In
module language it is the global relation shared by every vector, whereas
\(\operatorname{Ann}(v)\) records only the relation on one cyclic subspace.

\begin{lemma}
If \(V\) is finite-dimensional, then
\(\operatorname{Ann}_{F[x]}(V)=\{p\in F[x]:p(T)=0\}\) is nonzero.
\end{lemma}
\begin{proof}
The proof of Proposition~\ref{prop:operator-module-torsion} constructs, from
an \(F\)-basis \(v_1,\ldots,v_n\), nonzero polynomials \(f_i\) that
annihilate the individual basis vectors.  Their product annihilates every
basis vector and hence all of \(V\).
\end{proof}

Since \(F[x]\) is a PID, this annihilator has a unique monic generator,
called the \emph{minimal polynomial} \(m_T\):
\[
\operatorname{Ann}_{F[x]}(V_T)=(m_T).
\]
It is the monic polynomial of least degree satisfying \(m_T(T)=0\), and
every polynomial annihilating \(T\) is a multiple of it.

For \(v\in V\), write \(\operatorname{Ann}_{F[x]}(v)=(d_v)\) with \(d_v\)
monic. Since every global annihilator kills \(v\),
\((m_T)\subseteq(d_v)\), equivalently \(d_v\mid m_T\). If
\(v_1,\ldots,v_n\) is an \(F\)-basis and
\(\operatorname{Ann}(v_i)=(d_i)\), then a polynomial kills \(V\) exactly
when it kills every basis vector. Therefore
\[
\operatorname{Ann}(V_T)=\bigcap_{i=1}^n\operatorname{Ann}(v_i),
\qquad
m_T=\operatorname{lcm}(d_1,\ldots,d_n),
\]
with monic representatives. The minimal polynomial is thus the least common
polynomial relation satisfied simultaneously in every direction of \(V\).

\begin{proposition}[Eigenvalues and the minimal polynomial]
\label{prop:eigenvalue-minimal-polynomial}
For \(\lambda\in F\),
\[
\lambda\text{ is an eigenvalue of }T
\quad\Longleftrightarrow\quad m_T(\lambda)=0
\quad\Longleftrightarrow\quad (x-\lambda)\mid m_T.
\]
\end{proposition}
\begin{proof}
If \(Tv=\lambda v\) for some \(v\ne0\), then
\(0=m_T(T)v=m_T(\lambda)v\), so \(m_T(\lambda)=0\). The factor theorem gives
the second equivalence in one direction. Conversely, write
\(m_T=(x-\lambda)q\). If \(T-\lambda I\) were invertible, then
\[
0=m_T(T)=(T-\lambda I)q(T)
\]
would imply \(q(T)=0\), contradicting the minimal degree of \(m_T\). Thus
\(T-\lambda I\) is singular and \(\lambda\) is an eigenvalue.
\end{proof}

\begin{proposition}[Polynomial invariants of a direct sum]
\label{prop:direct-sum-polynomials}
Let \(T_i\in\mathcal L(V_i)\), \(i=1,2\), where the \(V_i\) are
finite-dimensional. Then
\[
m_{T_1\oplus T_2}
=\operatorname{lcm}(m_{T_1},m_{T_2}),
\qquad
\chi_{T_1\oplus T_2}=\chi_{T_1}\chi_{T_2},
\]
where the least common multiple is chosen monic.
\end{proposition}
\begin{proof}
A polynomial annihilates \(T_1\oplus T_2\) exactly when it annihilates both
summands.  Hence
\[
(m_{T_1\oplus T_2})
=(m_{T_1})\cap(m_{T_2}),
\]
whose monic generator in the PID \(F[x]\) is the least common multiple.
This also follows directly from
\(p(T_1\oplus T_2)=p(T_1)\oplus p(T_2)\). Choose bases and write
\(A_i=[T_i]\). In the concatenated basis,
\([T_1\oplus T_2]=[T_1]\oplus[T_2]\). Therefore
Chapter~5's block-determinant formula gives
\[
\det(xI-(A_1\oplus A_2))
=\det(xI-A_1)\det(xI-A_2).
\]
\end{proof}

\begin{theorem}[Cayley--Hamilton]\label{thm:cayley-hamilton}
Every linear operator satisfies its characteristic polynomial:
\[\chi_T(T)=0.\]
\end{theorem}
\begin{proofidea}
The matrix \(TI-A\) records the relations describing \(T\) on a basis.
Multiplying this relation matrix by its adjugate collapses it to its
determinant, forcing \(\chi_T(T)\) to kill every basis vector.
\end{proofidea}
\begin{proof}
Choose a basis \(\beta=(v_1,\ldots,v_n)\), write
\(A=(a_{ij})=[T]_\beta\), and regard the basis as the row vector
\(\mathbf v=(v_1,\ldots,v_n)\). For a matrix of operators \(Q=(q_{ij})\),
the \(j\)-th entry of \(\mathbf vQ\) means
\(\sum_i q_{ij}(v_i)\). Since column \(j\) of \(A\) gives the coordinates of
\(Tv_j\), the coordinate relations are exactly
\begin{equation}\label{eq:cayley-hamilton-relation-matrix}
\mathbf v(TI_n-A)=0.
\end{equation}

Every entry of \(TI_n-A\) lies in the subring
\[
F[T]=\{p(T):p\in F[x]\}\subseteq\operatorname{End}_F(V).
\]
Here each scalar entry \(a_{ij}\in F\) is identified with the scalar operator
\(a_{ij}I_V\in F[T]\), so every entry of \(TI_n-A\) genuinely belongs to
this one coefficient ring.
This ring is commutative because all its elements are polynomials in the same
operator. Determinants and adjugates therefore obey their usual identities
over \(F[T]\):
\[
(TI_n-A)\operatorname{adj}(TI_n-A)
=\det(TI_n-A)I_n.
\]
Multiplying \eqref{eq:cayley-hamilton-relation-matrix} on the right gives
\[
0=\mathbf v(TI_n-A)\operatorname{adj}(TI_n-A)
=\mathbf v\det(TI_n-A).
\]
Thus \(\det(TI_n-A)\) kills every basis vector.

It remains to identify this determinant. By definition,
\(\chi_T(x)=\det(xI_n-A)\). The evaluation homomorphism
\[
F[x]\longrightarrow F[T],\qquad p(x)\longmapsto p(T),
\]
takes \(\det(xI_n-A)\) to \(\det(TI_n-A)\), because a determinant is a
finite sum of products of entries. Hence
\[
\det(TI_n-A)=\chi_T(T).
\]
This operator kills every basis vector and therefore all of \(V\), proving
\(\chi_T(T)=0\).
\end{proof}

\begin{corollary}[Minimal polynomial divides characteristic polynomial]
\label{cor:minimal-divides-characteristic}
For every endomorphism of a finite-dimensional vector space,
\[
m_T\mid\chi_T,
\qquad \deg m_T\leq\dim V.
\]
\end{corollary}
\begin{proof}
Cayley--Hamilton says \(\chi_T\in\operatorname{Ann}(V_T)=(m_T)\), which is
exactly the divisibility assertion. Since \(\deg\chi_T=\dim V\), the degree
bound follows. Equivalently, polynomial division
\(\chi_T=qm_T+r\), \(\deg r<\deg m_T\), and evaluation at \(T\) give
\(r(T)=0\); minimality forces \(r=0\).
\end{proof}

If \(\chi_T\) splits over \(F\), then
Corollary~\ref{cor:minimal-divides-characteristic} shows that \(m_T\), as a
divisor of a split polynomial, also splits into linear factors. Every
irreducible \(p_i\) occurring in the operator module divides its annihilator
generator \(m_T\), so \(p_i=x-\lambda_i\) for some \(\lambda_i\in F\).
Proposition~\ref{prop:eigenvalue-minimal-polynomial} identifies the distinct
\(\lambda_i\) as the eigenvalues of \(T\), and
Proposition~\ref{prop:primary-generalized-eigenspace} turns the primary
decomposition from Section~\ref{sec:linear-operators-fx-modules} into
\[
V=G_{\lambda_1}(T)\oplus\cdots\oplus G_{\lambda_s}(T).
\]
This proves the generalized-eigenspace decomposition promised in Chapter~5,
using only results already established.

\section{Rational Canonical Form}
Rational canonical form exists over every field; it is the field-independent
classification of linear operators. Jordan form is available only when the
relevant irreducible polynomials split into linear factors over the base
field. Thus Jordan form is a split-polynomial refinement of rational form.

Recall from Section~\ref{sec:linear-operators-fx-modules} that the companion
matrix \(C_d\) is multiplication by \(x\) on the cyclic module
\(F[x]/(d)\) in the basis \(1,x,\ldots,x^{n-1}\). It is not an artificial
matrix recipe: it is the coordinate form of one generator subject to the
single polynomial relation \(d(x)=0\).

\begin{theorem}[Characteristic polynomial of a companion matrix]
\label{thm:companion-characteristic}
For a monic polynomial
\[
d(x)=x^n+a_{n-1}x^{n-1}+\cdots+a_0,
\]
the companion matrix \(C_d\) has characteristic polynomial
\[
\chi_{C_d}(x)=d(x).
\]
\end{theorem}
\begin{proof}
Write \(D_n(x)=\det(xI-C_d)\). Explicitly,
\[
\begin{pmatrix}
x&0&\cdots&0&a_0\\
-1&x&\cdots&0&a_1\\
0&-1&\ddots&0&a_2\\
\vdots&\ddots&\ddots&x&\vdots\\
0&\cdots&0&-1&x+a_{n-1}
\end{pmatrix}.
\]
Expand along the first row. Only the entries \(x\) and \(a_0\) contribute.
Deleting the first row and column leaves the companion-type determinant for
\[
x^{n-1}+a_{n-1}x^{n-2}+\cdots+a_2x+a_1;
\]
call it \(D_{n-1}(x)\). Hence the first contribution is
\(xD_{n-1}(x)\). The minor of \(a_0\) is triangular with diagonal entries
\(-1\), so its determinant is \((-1)^{n-1}\). Its cofactor sign is
\((-1)^{1+n}\), and
\[
(-1)^{1+n}(-1)^{n-1}=1.
\]
Thus
\[
D_n(x)=xD_{n-1}(x)+a_0.
\]
The case \(n=1\) is \(D_1(x)=x+a_0\). Induction now gives
\[
D_n(x)=x^n+a_{n-1}x^{n-1}+\cdots+a_1x+a_0=d(x).
\]
\end{proof}

The minimal polynomial of \(C_d\) is also \(d\), independently of the
determinant calculation. The matrix acts as multiplication by \(x\) on
\(F[x]/(d)\), whose class of \(1\) is cyclic. We have \(d(C_d)(1)=0\), while
no nonzero polynomial of degree less than \(n\) annihilates \(1\), because
such a polynomial cannot belong to the ideal \((d)\). Because the class of
\(1\) generates the entire \(F[x]\)-module, a polynomial annihilates the
whole module if and only if it annihilates that class. Therefore
\(m_{C_d}=d\).

\begin{lemma}[Characteristic polynomials and invariant subspaces]
\label{lem:invariant-subspace-characteristic-divisibility}
If \(W\leq V\) is \(T\)-invariant, then
\[
\chi_{T|_W}(x)\mid\chi_T(x).
\]
\end{lemma}
\begin{proof}
Choose a basis of \(W\) and extend it to a basis of \(V\). In this adapted
basis the matrix of \(T\) has block form
\[
\begin{pmatrix}[T|_W]&*\\0&C\end{pmatrix}.
\]
The block-triangular determinant formula,
Theorem~\ref{thm:block-triangular-determinant}, gives
\[
\chi_T(x)=\chi_{T|_W}(x)\det(xI-C).
\]
\end{proof}

\begin{corollary}[Polynomial invariants of a cyclic subspace]
\label{cor:cyclic-subspace-polynomials}
Use the notation of Proposition~\ref{prop:cyclic-subspace-annihilator}. Then
\[
[T|_W]_\beta=C_d,
\qquad
m_{T|_W}=\chi_{T|_W}=d,
\qquad
d\mid\chi_T.
\]
If \(v\) is cyclic for all of \(V\), so \(W=V\) and \(n=\dim V\), then
\[
\beta=(v,Tv,\ldots,T^{n-1}v),\qquad
[T]_\beta=C_d,\qquad
m_T=\chi_T=d.
\]
\end{corollary}
\begin{proof}
The proposition gives the natural orbit basis and
\([T|_W]_\beta=C_d\). The direct companion determinant computation gives
\(\chi_{T|_W}=d\). Since \(v\) generates \(W\), a polynomial annihilates
\(T|_W\) on all of \(W\) exactly when it annihilates \(v\); hence
\(\operatorname{Ann}_{F[x]}(W)=\operatorname{Ann}_{F[x]}(v)=(d)\) and
\(m_{T|_W}=d\). Finally, \(W\) is \(T\)-invariant, so
Lemma~\ref{lem:invariant-subspace-characteristic-divisibility} gives
\(d=\chi_{T|_W}\mid\chi_T\). Taking \(W=V\) yields the last display.
\end{proof}

\Needspace{10\baselineskip}
\begin{remark}[Alternative proof: cyclic invariant subspaces]
Here is a classical linear-algebra proof of Cayley--Hamilton. Fix
\(0\ne v\in V\) and put
\[
W=F[x]v.
\]
Corollary~\ref{cor:cyclic-subspace-polynomials} gives both
\[
\chi_{T|_W}(T)v=0
\qquad\text{and}\qquad
\chi_{T|_W}\mid\chi_T.
\]
Write
\(\chi_T=q\chi_{T|_W}\). Hence
\[
\chi_T(T)v=q(T)\chi_{T|_W}(T)v=0.
\]
The zero vector is automatic, and \(v\) was otherwise arbitrary, so
\(\chi_T(T)=0\).
\end{remark}

\begin{theorem}[Rational canonical form]\label{thm:rational-canonical-form}
Let \(T\) be an endomorphism of a finite-dimensional vector space over an
arbitrary field \(F\).  There are unique monic nonconstant polynomials
\[
d_1\mid d_2\mid\cdots\mid d_r
\]
such that, in a suitable basis,
\[
[T]=C_{d_1}\oplus\cdots\oplus C_{d_r}.
\]
The invariant factors \(d_i\) classify \(T\) up to similarity.  Moreover,
\[
m_T=d_r,
\qquad
\chi_T=d_1\cdots d_r.
\]
\end{theorem}
\begin{proof}
Proposition~\ref{prop:operator-module-torsion} makes \(V_T\) a finitely
generated torsion module over the PID \(F[x]\).  The structure theorem gives
its unique invariant factors and a decomposition
\[
V_T\cong\bigoplus_{i=1}^rF[x]/(d_i).
\]
The cyclic-subspace proposition turns multiplication by \(x\) on each
summand into \(C_{d_i}\).  The intertwiner proposition says that isomorphic
operator modules are exactly similar operators, proving both existence and
classification.

The annihilator of the direct sum is generated by the least common multiple
of the \(d_i\), which is \(d_r\) because they form a divisibility chain.
Thus \(m_T=d_r\).  The direct-sum characteristic-polynomial formula and
Theorem~\ref{thm:companion-characteristic} give
\(\chi_T=d_1\cdots d_r\).
\end{proof}

\begin{corollary}[Irreducible factors of the polynomial invariants]
\label{cor:minimal-characteristic-factors}
The polynomials \(m_T\) and \(\chi_T\) have exactly the same distinct monic
irreducible factors. Consequently, \(m_T\) splits over \(F\) if and only if
\(\chi_T\) splits over \(F\).
\end{corollary}
\begin{proof}
In rational canonical form, \(m_T=d_r\) and
\(\chi_T=d_1\cdots d_r\). Every monic irreducible factor of \(\chi_T\)
occurs in some \(d_i\); since \(d_i\mid d_r=m_T\), that irreducible factor
divides \(m_T\). Conversely, every monic irreducible factor of
\(m_T=d_r\) occurs in the product \(\chi_T=d_1\cdots d_r\).
\end{proof}

\begin{definition}
A vector \(v\in V\) is a \emph{cyclic vector} for \(T\) if
\(V=F[x]v\). The operator \(T\) is \emph{cyclic} if it has a cyclic vector.
\end{definition}

\begin{corollary}[Cyclic-operator criterion]
\label{cor:cyclic-operator-criterion}
For an endomorphism of a nonzero finite-dimensional vector space, the
following are equivalent:
\begin{enumerate}
\item \(T\) is cyclic;
\item its rational canonical form consists of one companion block;
\item \(m_T=\chi_T\).
\end{enumerate}
\end{corollary}
\begin{proof}
If \(V=F[x]v\), then
\(V_T\cong F[x]/\operatorname{Ann}(v)\), so its invariant-factor
decomposition has one summand and its rational form has one companion block.
Conversely, the class of \(1\) generates the module belonging to one
companion block, so such an operator is cyclic. A single block has both
minimal and characteristic polynomial equal to its defining polynomial.

Finally, if the rational form has invariant factors
\(d_1\mid\cdots\mid d_r\), then
\(m_T=d_r\) and \(\chi_T=d_1\cdots d_r\). All \(d_i\) are nonconstant, so
equality of these monic polynomials forces
\(\sum_i\deg d_i=\deg d_r\), hence \(r=1\).
\end{proof}

\begin{remark}[Cayley--Hamilton from the PID structure theorem]
The rational canonical form gives a second structural derivation. If
\(d_1\mid\cdots\mid d_r\) are the invariant factors, then
\[
V_T\cong\bigoplus_iF[x]/(d_i),\qquad
m_T=d_r,\qquad
\chi_T=d_1\cdots d_r.
\]
The last formula uses the direct companion determinant computation and the
block-diagonal determinant formula. Since \(d_r\mid d_1\cdots d_r\), we have
\(m_T\mid\chi_T\). Therefore \(m_T(T)=0\) implies
\(\chi_T(T)=0\). This argument is noncircular: the companion characteristic
polynomial was proved directly by cofactor expansion, without invoking
Cayley--Hamilton.
\end{remark}

\begin{remark}[Three views of Cayley--Hamilton]
The adjugate proof is a matrix identity from determinant algebra. The cyclic
invariant-subspace proof uses the characteristic polynomial of each local
cyclic restriction and its divisibility into \(\chi_T\). The PID proof says structurally that
\(\chi_T\) is a multiple of the annihilator generator \(m_T\):
\[
\begin{array}{c}
\text{adjugate identity}\longrightarrow\text{matrix proof},\\
\text{cyclic invariant subspaces}\longrightarrow\text{local restrictions and divisibility},\\
\text{PID structure theorem}\longrightarrow\text{module-theoretic proof}.
\end{array}
\]
Each proof exposes a different reason for the same theorem.
\end{remark}

\begin{example}[Rational form without real Jordan form]
Over \(\mathbb R\), let
\[
A=\begin{pmatrix}0&-1\\1&0\end{pmatrix}.
\]
Then \(\chi_A(x)=x^2+1\), and \(A^2=-I\) while no linear polynomial
annihilates \(A\), so \(m_A(x)=x^2+1\). Moreover,
\[
e_1=\begin{pmatrix}1\\0\end{pmatrix},
\qquad Ae_1=\begin{pmatrix}0\\1\end{pmatrix}
\]
form a basis, so \(e_1\) is cyclic and the operator module is explicitly
\(\mathbb R[x]/(x^2+1)\). Its rational canonical form is the companion
matrix of \(x^2+1\), which is exactly \(A\) in the convention above. There
is no real Jordan form because \(A\) has no real eigenvalues. Over
\(\mathbb C\), the factorization \(x^2+1=(x-i)(x+i)\) makes the operator
diagonalizable with eigenvalues \(i\) and \(-i\).
\end{example}

\begin{example}[From invariant factors to rational canonical form]
Suppose an \(F[x]\)-module has invariant factors
\[d_1=x-1,\qquad d_2=(x-1)(x^2+1).\]
Then it is the direct sum \(F[x]/(d_1)\oplus F[x]/(d_2)\). The first
summand contributes the \(1\times1\) companion matrix \((1)\), and the
second contributes the companion matrix of
\(x^3-x^2+x-1\). Their block diagonal sum is the rational canonical form.
At once we read
\[m_T=d_2=(x-1)(x^2+1),\qquad
\chi_T=d_1d_2=(x-1)^2(x^2+1).\]
The example illustrates the division of labor: invariant factors describe
the cyclic summands, while matrices make that description visible in a chosen
basis.
\end{example}

\begin{example}[One matrix through both canonical forms]
Over \(\mathbb C\), begin with the actual matrix
\[
A=\begin{pmatrix}1&1&0\\0&1&0\\0&0&-1\end{pmatrix}.
\]
Let \(T\) be its operator on \(\mathbb C^3\), and give the space its
\(\mathbb C[x]\)-module structure from
Section~\ref{sec:linear-operators-fx-modules}. On the first two standard
basis vectors, \((T-I)e_2=e_1\) and \((T-I)e_1=0\), so this plane is the
cyclic module \(\mathbb C[x]/((x-1)^2)\). The last basis vector gives
\(\mathbb C[x]/(x+1)\). Thus the elementary divisors are
\((x-1)^2\) and \(x+1\). They are coprime, so CRT combines them into the
single invariant factor
\[
d=(x-1)^2(x+1)=x^3-x^2-x+1.
\]
Consequently \(\chi_T=d\) and \(m_T=d\), the rational canonical form is
the companion matrix of \(d\),
\[
\begin{pmatrix}0&0&-1\\1&0&1\\0&1&1\end{pmatrix},
\]
and the Jordan form is \(J_2(1)\oplus(-1)\). The original matrix happens
already to be in Jordan form, but the calculation shows how its module
invariants produce both canonical descriptions rather than merely recognize
one by inspection.
\end{example}

\begin{example}[A noncanonical matrix through the classification]
Over \(\mathbb R\), consider
\[
A=\begin{pmatrix}2&0&1\\0&3&-1\\0&0&2\end{pmatrix}.
\]
Triangularity gives \(\chi_A(x)=(x-2)^2(x-3)\).  For \(N=A-2I\),
\[
\ker N=\operatorname{span}(e_1),
\qquad
\ker N^2=\operatorname{span}(e_1,e_2+e_3).
\]
Thus the generalized \(2\)-eigenspace has dimension \(2\), and
\[
N(e_2+e_3)=e_1,\qquad Ne_1=0
\]
is a Jordan chain of length \(2\).  The ordinary \(3\)-eigenspace is
\(\operatorname{span}(e_2)\).  Hence the elementary divisors are
\((x-2)^2\) and \(x-3\), so
\[
m_A=(x-2)^2(x-3)=\chi_A.
\]
Because the elementary divisors are coprime, they combine into the single
invariant factor
\[
d=(x-2)^2(x-3)=x^3-7x^2+16x-12.
\]
The rational canonical form is
\[
C_d=\begin{pmatrix}0&0&12\\1&0&-16\\0&1&7\end{pmatrix},
\]
while the Jordan canonical form is \(J_2(2)\oplus(3)\).

The Jordan basis \(\beta=(e_1,e_2+e_3,e_2)\) gives an explicit change of
basis.  With
\[
P=\begin{pmatrix}1&0&0\\0&1&1\\0&1&0\end{pmatrix}
\]
whose columns are the vectors of \(\beta\), one has
\[
P^{-1}AP=J_2(2)\oplus(3).
\]
This example starts from a matrix in neither canonical form and uses
generalized eigenspaces, kernel growth, a Jordan chain, elementary divisors,
and invariant factors to reach both forms.
\end{example}

\section{Jordan Canonical Form}
When the relevant polynomials split, elementary divisors refine the rational
blocks into the familiar Jordan blocks. This is the structural completion of
the diagonalization discussion in Chapter~5. The two canonical forms mirror
the two forms of the PID structure theorem:
\[
\begin{array}{ccl}
\text{invariant factors}&\longleftrightarrow&\text{rational companion blocks},\\
\text{elementary divisors}&\longleftrightarrow&\text{Jordan blocks}.
\end{array}
\]
Accordingly, Theorem~\ref{thm:uniqueness-invariant-factor} gives uniqueness
of rational canonical form, while
Theorem~\ref{thm:uniqueness-elementary-divisor} gives uniqueness of the
Jordan blocks.

\begin{theorem}[Jordan canonical form]\label{thm:jordan-canonical-form}
Let \(T\) be an endomorphism of a finite-dimensional \(F\)-vector space and
suppose that \(\chi_T\) splits over \(F\).  Then \(V\) has a basis in which
\[
[T]=J_{k_1}(\lambda_1)\oplus\cdots\oplus J_{k_s}(\lambda_s).
\]
The blocks correspond to the elementary divisors
\[
(x-\lambda_1)^{k_1},\ldots,(x-\lambda_s)^{k_s},
\]
and their multiset is unique up to ordering.  Equivalently, \(V\) has a
basis obtained by concatenating Jordan chains.
\end{theorem}
\begin{proof}
When \(\chi_T\) splits, every irreducible factor in the elementary-divisor
decomposition of \(V_T\) is \(x-\lambda\).  Thus the PID structure theorem
gives a direct sum of cyclic modules
\[
F[x]/((x-\lambda)^k).
\]
Proposition~\ref{prop:jordan-chain-cyclic-module} supplies a Jordan-chain
basis and the block \(J_k(\lambda)\) on each summand.  Uniqueness of the
elementary divisors makes the multiset of blocks unique immediately; no
separate linear-algebraic uniqueness proof is needed.
\end{proof}

\begin{corollary}[Diagonalizability criterion]
\label{cor:diagonalizable-minimal-polynomial}
An endomorphism \(T\) is diagonalizable over \(F\) if and only if \(m_T\) is
a product of distinct linear factors over \(F\).
\end{corollary}
\begin{proof}
In Jordan form, diagonalizability says exactly that every block has size
\(1\). The elementary divisor belonging to a block \(J_k(\lambda)\) is
\((x-\lambda)^k\), and \(m_T\) is the least common multiple of all elementary
divisors. Thus all block sizes are \(1\) exactly when \(m_T\) splits into
distinct linear factors; for the reverse implication,
Corollary~\ref{cor:minimal-characteristic-factors} first ensures that
\(\chi_T\) splits, so Jordan form exists. The splitting condition is
essential: real rotation
of the plane has no real eigenbasis.
\end{proof}

\begin{corollary}[Largest blocks and kernel growth]
\label{cor:jordan-kernel-growth}
Assume that \(\chi_T\) splits over \(F\). For each eigenvalue \(\lambda\),
the exponent of \(x-\lambda\) in \(m_T\)
is the size of the largest Jordan block for \(\lambda\). Moreover, put
\[
K_j=\ker(T-\lambda I)^j,
\qquad K_0=0,
\qquad c_j=\dim K_j-\dim K_{j-1}.
\]
Then
\[
c_j=\#\{\text{Jordan blocks for \(\lambda\) of size at least \(j\)}\},
\]
and therefore
\[
c_j-c_{j+1}
=\#\{\text{Jordan blocks for \(\lambda\) of size exactly \(j\)}\}.
\]
Finally,
\[
\dim G_\lambda(T)=
\text{the algebraic multiplicity of \(\lambda\)}.
\]
\end{corollary}
\begin{proof}
The least common multiple of the elementary divisors contains
\((x-\lambda)^k\) to the largest exponent occurring among the
\(\lambda\)-blocks, proving the first assertion. On a block of size \(k\)
for \(\lambda\), the kernel of the \(j\)-th power has dimension
\(\min(j,k)\), so its contribution to \(c_j\) is \(1\) exactly when
\(k\geq j\). On a block for a different eigenvalue, \(T-\lambda I\) is
invertible and the contribution is zero. Summing over blocks proves both
counting formulas.

The generalized eigenspace is the direct sum of all \(\lambda\)-blocks, so
its dimension is their total size. The same total is the exponent of
\(x-\lambda\) in \(\chi_T\), hence the algebraic multiplicity.
\end{proof}

Chapter~5 established these kernel-growth formulas conditionally from a
displayed Jordan decomposition. The Jordan theorem supplies the required
decomposition for every split characteristic polynomial, so the bookkeeping
is now unconditional under precisely that hypothesis.

\begin{example}[From elementary divisors to Jordan form]
Over \(\mathbb C\), suppose the elementary divisors are
\[(x-2)^3,\qquad (x-2),\qquad (x+1)^2.\]
There are therefore two Jordan blocks for eigenvalue \(2\), of sizes \(3\)
and \(1\), and one block for eigenvalue \(-1\), of size \(2\). The Jordan
form is
\[J_3(2)\oplus J_1(2)\oplus J_2(-1).\]
Its characteristic polynomial is \((x-2)^4(x+1)^2\), whereas its minimal
polynomial is \((x-2)^3(x+1)^2\): the largest block for each eigenvalue
determines the corresponding exponent in the minimal polynomial.
\end{example}

\begin{remark}[How the polynomial invariants fit together]
If \(d_1\mid\cdots\mid d_r\) are invariant factors, then
\(\chi_T=d_1\cdots d_r\) and \(m_T=d_r\). Factoring the \(d_i\)'s and
using CRT yields elementary divisors. When these factors are powers of linear
polynomials, their exponents are exactly the Jordan-block sizes. Thus the
characteristic polynomial records total algebraic multiplicities, the minimal
polynomial records the largest required exponents, invariant factors produce
rational canonical form over every field, and elementary divisors produce
Jordan form when splitting is available.
\[
\begin{gathered}
\text{PID structure theorem}\\
\big\downarrow\\
\text{invariant factors}\\
\big\downarrow\\
\text{rational canonical form over every field},\\[1.5ex]
\text{factor into prime powers}\\
\big\downarrow\\
\text{elementary divisors}\\
\big\downarrow\\
\text{Jordan blocks when }p=x-\lambda.
\end{gathered}
\]
Rational canonical form is the field-independent classification.  Jordan
form is its split-polynomial refinement.
\end{remark}

\section*{Exercises}
\begin{exercise}[Basic] Classify
\[
\mathbb Z/8\mathbb Z\oplus\mathbb Z/12\mathbb Z
\]
in invariant-factor and elementary-divisor form.
\end{exercise}
\begin{exercise}[Core] Convert
\[
\mathbb Z/4\oplus\mathbb Z/8\oplus\mathbb Z/9\oplus\mathbb Z/3
\]
from elementary-divisor form to invariant-factor form, and convert
\(\mathbb Z/12\oplus\mathbb Z/60\) in the opposite direction.
\end{exercise}
\begin{exercise}[Core] List the abelian groups of order \(360=2^3 3^2 5\)
up to isomorphism by choosing partitions prime by prime.  Check the order of
each group from its invariant factors.\end{exercise}
\begin{exercise}[Core] Find the Smith normal form of
\(\begin{pmatrix}2&4&6\\4&10&14\end{pmatrix}\) over \(\mathbb Z\), regarded
as a homomorphism \(\mathbb Z^3\to\mathbb Z^2\), and determine its
cokernel.\end{exercise}
\begin{exercise}[Core] For \(T(x,y)=(-y,x)\) on \(\mathbb R^2\), find the
minimal polynomial and explain why no real Jordan form exists.\end{exercise}
\begin{exercise}[Core] Let \(T:V\to V\), \(S:W\to W\), and let
\(\phi:V\to W\) be \(F\)-linear.  Verify directly for
\(p(x)=2x^3-x+1\) that \(\phi T=S\phi\) implies
\(\phi p(T)=p(S)\phi\), and then prove the assertion for every
\(p\in F[x]\).\end{exercise}
\begin{exercise}[Core] For a vector \(v\), suppose
\(v,Tv,T^2v,T^3v\) are independent and
\[
T^4v=2v-Tv+3T^3v.
\]
Find \(\operatorname{Ann}(v)\) and the companion matrix of \(T\) on
\(F[x]v\).\end{exercise}
\begin{exercise}[Core] Prove the minimal-polynomial formula for a direct sum
of three operators by applying Proposition~\ref{prop:direct-sum-polynomials}
twice.  Give an example showing why the answer is an lcm rather than a
product.\end{exercise}
\begin{exercise}[Core] For
\[
A=\begin{pmatrix}1&1\\0&1\end{pmatrix}\in M_2(F),
\]
compute the Smith normal form of \(xI-A\) over \(F[x]\), recover the invariant
factors of \(V_A\), and determine the rational canonical form.\end{exercise}
\begin{exercise}[Further] Show that the invariant factors of a diagonalizable
operator are square-free.\end{exercise}
\begin{exercise}[Further] First suppose an operator has
\[
\dim\ker(T-\lambda I)^j=2,4,5,6,6,\ldots,
\]
and recover the sizes of all its Jordan blocks for \(\lambda\).  In a
separate example with different block data, construct explicit chains for
the matrix
\[
\begin{pmatrix}
\lambda&1&0&0\\0&\lambda&0&0\\
0&0&\lambda&1\\0&0&0&\lambda
\end{pmatrix}.
\]\end{exercise}
\begin{exercise}[Further] Decide whether the matrices
\[
\begin{pmatrix}1&1&0\\0&1&0\\0&0&2\end{pmatrix},
\qquad
\begin{pmatrix}1&0&0\\0&1&1\\0&0&2\end{pmatrix}
\]
are similar.  Justify the answer using elementary divisors rather than an
attempt to solve directly for a change-of-basis matrix.\end{exercise}
\begin{exercise}[Further, Challenging] Determine the rational and Jordan
canonical forms of an operator whose elementary divisors are
\((x-1)^3,(x-1),(x+2)^2\).\end{exercise}
\begin{exercise}[Looking Ahead] For a Jordan matrix, read off
\(\chi_T\), \(m_T\), and the dimensions of all generalized eigenspaces.
For a rational canonical matrix, read off \(\chi_T\) and \(m_T\) from its
companion blocks and explain how factorization over a splitting field predicts
the Jordan blocks.\end{exercise}
